\chapter{Architecture}
\label{chapter:architecture}

This chapter describes the current and proposed architectures for managing vulnerability data, emphasizing modularity, scalability, and efficient interaction between system components. Building on the multi-dimensional classification and remediation model outlined in chapter \ref{chapter:multidimensional-vulnerability-classification}, the proposed extended architecture translates these concepts into a scalable and practical system design.

\section{Current Architecture}
\label{sec:current-architecture}

The current solution integrates data from the \ac{OSV} database into the internal system for effective vulnerability management. The process ensures efficient synchronization, processing, and storage of vulnerability advisories, aligning them with components defined in software bill of materials (\ac{SBOM}s) to identify affected packages and components with vulnerabilities.

\begin{itemize}
    \item \textbf{Data Synchronization:} The system initiates synchronization by checking for the availability of the vulnerability data archive (\ac{OSV}) in a cloud-based storage system. This includes mechanisms to verify the freshness of the data (e.g., through metadata such as ETags) to avoid redundant processing. Only updated or newly added archives are retrieved and temporarily stored for processing. This process happens once a day.

    \item \textbf{Data Extraction and Processing:} Retrieved archives are decompressed, and individual advisories are parsed and processed. Each advisory is evaluated for updates since the last synchronization and is enriched with relevant details, such as identifiers, severity metrics, and associated packages. 
    \begin{itemize}
        \item \textbf{Change Detection:} Only advisories that have been modified since the previous synchronization are processed further.
        \item \textbf{Mapping and Enrichment:} Relevant details, including affected components, vulnerability types, and risk metrics (e.g., \ac{CVSS}), are extracted and structured for database integration.
    \end{itemize}

    \item \textbf{Component Alignment:} Processed advisories are matched against the components listed in the \ac{SBOM}s. This step identifies specific packages and software elements within the \ac{SBOM} that are directly affected by vulnerabilities, enabling precise vulnerability reporting at the component level.

    \item \textbf{Database Update:} The structured advisories, enriched with mappings to specific components, are stored in the internal database. Synchronization metadata, such as timestamps and version identifiers, is updated to reflect the latest state of external data sources. When an \ac{API} request is triggered, the relevant data is sent to the frontend, ensuring the user interface is always in sync with the most recent vulnerability information.
\end{itemize}

\section{Extended Architecture}
\label{sec:extended-architecture}

This design centers on a streamlined approach to vulnerability information management. It integrates multiple data sources (e.g., vulnerability databases, scoring providers) and applies caching to reduce unnecessary network requests. Its purpose is to unify critical metrics (such as \ac{CVSS}, \ac{EPSS}) into a consolidated view, as described in the multi-dimensional classification model (see section \ref{sec:multi-dimensional-classification}), while also supporting role-specific actions based on the remediation model outlined in section \ref{sec:remediation_recommendation}.

\subsection{Backend Responsibilities}
\label{subsec:backend-responsibilities}

The backend is responsible for data retrieval, caching, and score computation, following the scoring algorithm described in section \ref{sec:algorithm-classifications}:
\begin{itemize}
    \item \textbf{Multi-Source Lookup:} For each vulnerability, the system first checks existing records in an internal repository to discover any \ac{CVE} identifier or \ac{CVSS} vector. If missing, it fetches those details from external sources (e.g., a global advisory service or an official database). The \ac{EPSS} metric is always retrieved externally to ensure it is up to date, as it is not included in the internal repository.
    \item \textbf{Caching and Time-Based Validation:} Before making an external request, the system verifies whether relevant data (e.g., \ac{CVSS} details) is still current. If it is, it serves from the cache. Otherwise, it triggers a fresh fetch, updates the cache, and recalculates any derived metrics (e.g., severity).
    \item \textbf{Role-Specific Recommendations:} Based on user input (e.g., developer or security advisor context), the system determines a recommended course of action. The backend solely stores the \ac{SSVC} recommendation, which might include immediate patching, hotfixes, or a scheduled approach, for later retrieval or review.
\end{itemize}

\subsection{Frontend Responsibilities}
\label{subsec:frontend-responsibilities}

The frontend is responsible for fetching data from the backend, displaying it to the user, and forwarding user inputs for storage. Key functionalities include:

\begin{itemize}
  \item \textbf{View Scores and Vectors:} A dedicated interface fetches data from the backend and displays \ac{CVE} identifiers, \ac{CVSS} vectors, \ac{CVSS} base scores, \ac{EPSS} values, overall severity, and missing data, as described in section \ref{subsec:ui-explaining-score}. Additionally, it is responsible for showing vulnerabilities ranked as described in section \ref{sec:rank-order-vulnerabilities}.
  \item \textbf{Perform Structured Assignments:} An interactive decision flow guides users (depending on their role) through a series of questions, ultimately generating a recommendation, as seen in the concept of section \ref{sec:remediation_recommendation}. 
  \item \textbf{Send Updates to the Backend:} Once the recommendation is finalized, it is transmitted to the backend for persistence.
\end{itemize}


\subsection{System Interaction Workflow}
\label{subsec:system-interaction}

This section describes the mechanisms for fetching and caching vulnerability data.

\begin{figure}[H]
\centering
\begin{tikzpicture}[
    scale=0.625,
    transform shape,
    node distance=3.75cm,
    every node/.style={font=\scriptsize},
    >=latex',
    thick
]

\tikzstyle{decision} = [
    diamond,
    draw,
    fill=blue!5,
    text width=6em,
    text centered,
    inner sep=0pt
]
\tikzstyle{block} = [
    rectangle,
    draw,
    fill=green!5,
    text width=6em,
    text centered,
    rounded corners,
    minimum height=2.5em
]
\tikzstyle{line} = [draw, -latex']

\node[block] (start) {Start API-Request};

\node[decision, below of=start] (vulndata) {Entity up to date\\ or missing?};
\path[line] (start) -- (vulndata);

\node[draw, fill, circle,
      below=14cm of start,
      xshift=4.54cm,
      label={[font=\small]below:{Return to caller}},
      minimum size=30pt,      
      inner sep=0pt]         
(returnToCaller) {};

\draw[line]
  (vulndata.east) -- ++(3,0) node[above]{up to date}
  -- (returnToCaller);

\node[block, left=3.5cm of vulndata] (create) {Create entity};
\path[line] (vulndata) -- node[above] {missing} (create);

\node[decision, below of=create] (cveid) {CVE-ID\\ found in cache?};
\path[line] (create) -- (cveid);

\node[decision, below of=cveid] (cvss) {CVSS-Details\\ found in cache?};

\node[block, right=4.5cm of cveid] (fetch) {fetchFromApi};

\node[decision, below of=fetch] (fetchOK) {Fetch\\ success?};

\node[block, right=4.5cm of fetchOK] (error) {Error};

\node[block, below of=cvss] (epss) {getEpssScore};

\node[block, below=1cm of fetchOK] (calculate) {calculate\\ severity score};

\node[block, below=1cm of calculate] (store) {save collected\\ data};

% cveid -> no -> fetch
\path[line] (cveid.east) -- ++(0.8,0) node[above]{no} |- (fetch);

% cveid -> yes -> cvss
\path[line] (cveid) -- node[left]{yes} (cvss);

% cvss -> no -> fetch
\path[line] (cvss.east) -- ++(0.8,0) node[right]{no} |- (fetch);

% cvss -> yes -> epss
\path[line] (cvss) -- node[left]{yes} (epss);

% fetch -> fetchOK
\path[line] (fetch) -- (fetchOK);

% fetchOK -> yes -> epss
\path[line] (fetchOK.south) -- (calculate.north) node[midway,left]{yes};

% fetchOK -> no -> error
\path[line] (fetchOK.east) -- ++(0.8,0) node[above]{no} 
    |- (error);

% error -> ReturnToCaller
\path[line] (error) |- (returnToCaller);

% epss -> calculate
\path[line] (epss.east) -- ++(2,0) |- (fetch.west);

% calculate -> store
\path[line] (calculate) -- (store);

% store -> ReturnToCaller
\path[line] (store) -- (returnToCaller);

\path[line] (vulndata.south) -- ++(0,-0.75) node[midway,right]{outdated} -- (cveid.north);

\end{tikzpicture}
\caption{The internal data-fetch and chaching flow for vulnerability data.}
\label{fig:data-fetch-flow}
\end{figure}

The data-fetch and caching flow, illustrated in figure \ref{fig:data-fetch-flow}, begins when the frontend initiates a request to the backend, which checks whether the targeted vulnerability entity is already up to date. If it is, the system returns immediately to the caller; if the entity is missing or outdated, the workflow creates a new record or updates the existing one. Next, the backend determines whether the \ac{CVE} identifier is present in its cache. Finding it in the cache leads to a further check of whether \ac{CVSS} details are also available. If not, the system calls \texttt{fetchFromApi} to retrieve the necessary information from an external source.

When the backend looks for \ac{CVSS} data and does not find it in the cache, it again invokes \texttt{fetchFromApi}. If these fetch operations are successful, the retrieved data forms the basis for calculating an overall severity score (by combining \ac{CVSS} and \ac{EPSS} metrics, as detailed in section \ref{sec:algorithm-classifications}). A failed retrieval, however, diverts the workflow to an \texttt{Error} state, where a 10.1 (see section \ref{subsec:handling-missing-data}) severity score (highlighted in the frontend in pink) is returned. In such cases, users can refer to a dedicated \emph{Score Explanation} feature on the frontend (see section \ref{subsec:ui-explaining-score}) to learn how to remedy the issue. Once the \ac{EPSS} data has been received, the backend finalizes its calculations and saves all new or updated information to an internal repository. The process concludes when the system returns control to the caller, providing fully updated vulnerability records.

\subsection{Rate Limits and Caching Benefits}
\label{subsec:rate-limits-caching}

External data sources such as the \ac{NIST} \ac{NVD}, the \ac{GAD}, and the \ac{EPSS} all impose rate limits to protect against misuse. In this architecture, \ac{CVE} data is fetched from the \ac{GAD}, \ac{CVSS} data from the \ac{NVD}, and \ac{EPSS} data from \texttt{first.org}. For example, \ac{NIST} limits requests to five per 30-second window without an API key, and fifty when using one.\footnote{\url{https://nvd.nist.gov/developers/start-here}} The \ac{GAD} allows unauthenticated users sixty requests per hour and offers 5.000 or more requests per hour for authenticated users, with higher ceilings for GitHub Enterprise organizations.\footnote{\url{https://docs.github.com/en/rest/overview/resources-in-the-rest-api\#rate-limiting}} Meanwhile, the \ac{EPSS} service enforces a threshold of 1.000 requests per hour without a token.\footnote{\url{https://api.first.org/\#Rate-Limit}} 

To cope with these constraints, the backend employs caching, thereby reducing repetitive lookups. It fetches data only when a record is missing or deemed outdated, which significantly cuts down on external calls and helps avoid surpassing rate limits. As a result, the system performs more efficiently and can handle errors gracefully by returning a score of 10.1 in the event of being throttled or encountering other fetch errors, which is displayed in the frontend with a pink highlight. Users then have the option to consult the score explanation interface to determine the next steps for resolving the issue or reattempting data retrieval.

\subsection{SSVC Process and Role-Specific Assignments}
\label{subsec:ssvc-process}

The \ac{SSVC} process involves tailoring vulnerability handling to each user's or team's specific context, following the stakeholder-specific prioritization framework described in section \ref{sec:algorithm-remediation}. This process is supported by the controller layer, which manages endpoints like the \ac{SSVC} recommendation endpoint. These endpoints facilitate the submission of role-specific decisions, which are then persisted and integrated with other metrics (e.g., \ac{CVSS}, \ac{EPSS}). Figure~\ref{fig:ssvc-interaction} illustrates this interaction.

\begin{figure}[H]
\centering
\resizebox{0.9\textwidth}{!}{%
\begin{tikzpicture}[>=stealth', node distance=2.5cm, auto, thick, font=\sffamily\small]

  % Nodes
  \node[draw, rectangle, rounded corners, align=center, minimum width=3cm, minimum height=1cm] (frontend) {User Interface};
  \node[draw, rectangle, rounded corners, right=4.0cm of frontend, minimum width=3.3cm, align=center] (backendController) {Controller Layer};
  \node[draw, rectangle, rounded corners, right=3.5cm of backendController, minimum width=3.3cm, align=center] (backendStorage) {Data Storage \\ \& Calculation};

  % Paths with slight curves to avoid overlap
  \draw[->] (frontend.east)
    to[out=0, in=180]
    node[above, yshift=0.2cm]{\scriptsize Finalize Role-Based Decision}
    (backendController.west);

  \draw[->] (backendController.east)
    to[out=0, in=180]
    node[above, yshift=0.2cm]{\scriptsize Persist and Merge Data}
    (backendStorage.west);

  \draw[->] (backendStorage.south west)
    to[out=230, in=-50]
    node[below, yshift=-0.2cm]{\scriptsize Return Updated Metrics}
    (backendController.south east);

  \draw[->] (backendController.south west)
    to[out=230, in=-50]
    node[below, yshift=-0.2cm]{\scriptsize Updated Values (CVSS, EPSS, etc.)}
    (frontend.south east);

\end{tikzpicture}
}
\caption{High-level interaction in the SSVC process.}
\label{fig:ssvc-interaction}
\end{figure}

\section{Design Principles}
\label{sec:design-principles}

\begin{itemize}
  \item \textbf{Modularity:} Separate retrieval, caching, scoring, and role-based recommendation processes, allowing each to be maintained independently.
  \item \textbf{Scalability:} Use time-based refresh thresholds and micro-caching to handle growing numbers of vulnerabilities and high API traffic.
  \item \textbf{Usability:} Provide a clear sequence for role-based input, ensuring that developers or security advisors receive context-specific guidance without being overwhelmed by unrelated details.
  \item \textbf{Consistency:} Apply uniform approaches to caching and scoring so that users have a predictable view of vulnerability metrics, regardless of the external source.
\end{itemize}

\section{Conclusion}
\label{sec:architecture-conclusion}
The outlined architecture consolidates key vulnerability metrics and merges them with role-based inputs to yield tailored recommendations. By systematically checking internal data first and fetching new information only when required, it minimizes redundant requests while maintaining data accuracy. Frontend interactions are kept straightforward through specialized dialogs and status views, ensuring users can easily view or update vulnerability details. This approach offers a strong foundation for integrating further data sources, automating compliance checks, and scaling to large application ecosystems.